\documentclass[11pt,a4paper]{article}

% ---------- packages ----------
\usepackage[utf8]{inputenc}
\usepackage[T1]{fontenc}
\usepackage{booktabs}
\usepackage{graphicx}
\usepackage{amsmath}
\usepackage{hyperref}
\usepackage[margin=1in]{geometry}
\usepackage{natbib}
\usepackage{xcolor}
\usepackage{multirow}

% ---------- placeholder command ----------
\newcommand{\PH}[1]{{\color{red}\textbf{#1}}}

% ---------- metadata ----------
\title{Evaluating LLM Claim Extraction Accuracy\\
for Autonomous Knowledge Graph Construction}

\author{
  Autonomous AI Researcher \\
  PicoClaw Research \\
  \texttt{research@picoclaw.dev}
}

\date{\today}

% ==========================================================================
\begin{document}
\maketitle

% ---------- abstract ----------
\begin{abstract}
Autonomous research systems depend on accurate extraction of factual claims
from scientific papers to build reliable knowledge graphs.  We evaluate three
large language models---Claude Sonnet 4, GPT-4o, and Claude Haiku 4---on
the task of extracting verifiable claims from machine-learning paper
abstracts.  Using a dual-annotator ground-truth protocol and embedding-based
claim matching (cosine similarity threshold 0.85), we measure per-model
precision, recall, F1, and hallucination rate.  Our best-performing model,
\PH{{{BEST\_MODEL}}}, achieves a precision of \PH{{{PRECISION\_BEST}}},
recall of \PH{{{RECALL\_BEST}}}, and F1 of \PH{{{F1\_BEST}}}, with a
hallucination rate of \PH{{{HALLUCINATION\_BEST}}}.  We test three
pre-registered hypotheses (precision $\geq 0.80$, recall $\geq 0.70$,
hallucination rate $\leq 0.10$) and report bootstrap 95\% confidence
intervals and pairwise McNemar tests.  Results indicate that
\PH{{{SUMMARY\_SENTENCE}}}.  These findings establish baseline quality
thresholds for claim extraction in downstream autonomous research pipelines.
\end{abstract}

% ==========================================================================
\section{Introduction}
\label{sec:introduction}

Building reliable knowledge graphs from scientific literature is a
prerequisite for autonomous AI research systems that must reason over an
evolving body of evidence~\citep{lu2024aiscientist}.  The first step in
such a pipeline is \emph{claim extraction}: given a paper (or its
abstract), produce a set of discrete, verifiable factual statements.

Errors at this stage propagate downstream: a spurious claim (hallucination)
can trigger wasted experiments, while a missed claim creates blind spots
in the literature review.  Quantifying these failure modes across frontier
language models is therefore essential before deploying an autonomous
research agent.

In this work we design a controlled experiment (Experiment~E01) that
measures the precision, recall, F1, and hallucination rate of three LLMs
on claim extraction from ML paper abstracts.  We contribute:

\begin{enumerate}
  \item A dual-annotator ground-truth corpus of claims extracted from
        \PH{{{N\_PAPERS}}} ArXiv abstracts, with reconciliation and
        inter-annotator agreement statistics.
  \item A systematic comparison of Claude Sonnet 4, GPT-4o, and
        Claude Haiku 4 on this corpus.
  \item Pre-registered hypothesis tests with bootstrap confidence intervals,
        McNemar pairwise comparisons, and effect-size estimates.
\end{enumerate}

% ==========================================================================
\section{Related Work}
\label{sec:related}

\paragraph{LLM-as-Judge.}
\citet{zheng2023judging} demonstrate that strong LLMs can serve as reliable
evaluators of open-ended text generation, achieving high agreement with
human raters.  We extend this paradigm to structured claim extraction,
where the evaluation target is a \emph{set} of propositions rather than a
single response.

\paragraph{AI Scientist.}
\citet{lu2024aiscientist} present an end-to-end autonomous research agent
capable of ideation, experimentation, and paper writing.  Their system
relies on accurate literature understanding but does not isolate claim
extraction as a measurable subtask.  Our experiment fills this gap by
providing quantitative baselines.

\paragraph{Information extraction from scientific papers.}
Prior work on scientific information extraction has focused on named-entity
recognition~\citep{beltagy2019scibert}, relation extraction, and citation
intent classification.  Claim-level extraction---identifying self-contained
factual assertions---is less studied, although recent work on fact
verification~\citep{wadden2020fact} addresses the complementary problem of
checking claims against evidence.

% ==========================================================================
\section{Methodology}
\label{sec:methodology}

\subsection{Paper Selection}
\label{sec:paper-selection}

We sample \PH{{{N\_PAPERS}}} abstracts from the ArXiv preprint server
using stratified sampling across \PH{{{N\_CATEGORIES}}} machine-learning
subcategories (cs.LG, cs.CL, cs.AI, stat.ML).  Papers are drawn from
the period \PH{{{DATE\_RANGE}}} and filtered to ensure a minimum abstract
length of 100 words.

\subsection{Ground Truth Annotation}
\label{sec:annotation}

Two independent annotators (LLM-based, using a separate model instance
from those under evaluation) extract claims from each abstract under two
annotation strategies:

\begin{description}
  \item[Conservative:] Only claims that are explicitly and unambiguously
    stated in the text.
  \item[Comprehensive:] All claims, including those reasonably implied by
    the text.
\end{description}

We reconcile the two annotation sets by taking the union and adjudicating
disagreements.  Inter-annotator agreement is measured with Cohen's
$\kappa$: \PH{{{COHENS\_KAPPA}}}.

\subsection{Models Tested}
\label{sec:models}

We evaluate three models accessed via API:

\begin{enumerate}
  \item \textbf{Claude Sonnet 4} (Anthropic, direct API).
  \item \textbf{GPT-4o} (OpenAI, direct API).
  \item \textbf{Claude Haiku 4} (Anthropic, via OpenRouter).
\end{enumerate}

Each model receives the same system prompt and paper abstract.  We perform
\PH{{{N\_RUNS}}} independent runs per model per paper to assess extraction
stability.

\subsection{Claim Matching}
\label{sec:matching}

Extracted claims are matched to ground-truth claims using embedding cosine
similarity.  We embed all claims with \PH{{{EMBEDDING\_MODEL}}} and define
a match when the cosine similarity exceeds a threshold of $\tau = 0.85$.
Each ground-truth claim is matched to at most one extracted claim
(bipartite matching).

\begin{itemize}
  \item \textbf{True positive (TP):} extracted claim matches a
        ground-truth claim ($\mathrm{sim} \geq 0.85$).
  \item \textbf{False positive (FP):} extracted claim has no match
        (potential hallucination).
  \item \textbf{False negative (FN):} ground-truth claim has no match.
\end{itemize}

\subsection{Statistical Tests}
\label{sec:statistics}

We report the following statistics:

\begin{itemize}
  \item \textbf{Bootstrap 95\% CIs} ($B = 10{,}000$) for precision,
        recall, and F1 per model.
  \item \textbf{McNemar's test} for pairwise comparison of model accuracy
        on matched claim pairs.
  \item \textbf{Cohen's $\kappa$} for inter-annotator agreement.
  \item \textbf{Binomial 95\% CI} for hallucination rate.
  \item \textbf{Cohen's $d$} for effect size between model pairs on F1.
\end{itemize}

% ==========================================================================
\section{Results}
\label{sec:results}

\subsection{Per-Model Performance}

Table~\ref{tab:per-model} summarises precision, recall, and F1 with 95\%
bootstrap confidence intervals for each model.

\begin{table}[ht]
\centering
\caption{Per-model claim extraction performance (95\% bootstrap CI).}
\label{tab:per-model}
\begin{tabular}{lccc}
\toprule
\textbf{Model} & \textbf{Precision} & \textbf{Recall} & \textbf{F1} \\
\midrule
Claude Sonnet 4
  & \PH{{{PRECISION\_SONNET}}} \scriptsize[\PH{{{CI\_PREC\_SONNET}}}]
  & \PH{{{RECALL\_SONNET}}}    \scriptsize[\PH{{{CI\_REC\_SONNET}}}]
  & \PH{{{F1\_SONNET}}}        \scriptsize[\PH{{{CI\_F1\_SONNET}}}] \\
GPT-4o
  & \PH{{{PRECISION\_GPT4O}}} \scriptsize[\PH{{{CI\_PREC\_GPT4O}}}]
  & \PH{{{RECALL\_GPT4O}}}    \scriptsize[\PH{{{CI\_REC\_GPT4O}}}]
  & \PH{{{F1\_GPT4O}}}        \scriptsize[\PH{{{CI\_F1\_GPT4O}}}] \\
Claude Haiku 4
  & \PH{{{PRECISION\_HAIKU}}} \scriptsize[\PH{{{CI\_PREC\_HAIKU}}}]
  & \PH{{{RECALL\_HAIKU}}}    \scriptsize[\PH{{{CI\_REC\_HAIKU}}}]
  & \PH{{{F1\_HAIKU}}}        \scriptsize[\PH{{{CI\_F1\_HAIKU}}}] \\
\bottomrule
\end{tabular}
\end{table}

\subsection{Pairwise Model Comparisons}

Table~\ref{tab:mcnemar} reports pairwise McNemar test results and Cohen's
$d$ effect sizes.

\begin{table}[ht]
\centering
\caption{Pairwise McNemar test ($\chi^2$, $p$-value) and Cohen's $d$.}
\label{tab:mcnemar}
\begin{tabular}{lccc}
\toprule
\textbf{Comparison} & $\chi^2$ & $p$ & Cohen's $d$ \\
\midrule
Sonnet vs.\ GPT-4o
  & \PH{{{CHI2\_SONNET\_GPT4O}}}
  & \PH{{{P\_SONNET\_GPT4O}}}
  & \PH{{{D\_SONNET\_GPT4O}}} \\
Sonnet vs.\ Haiku
  & \PH{{{CHI2\_SONNET\_HAIKU}}}
  & \PH{{{P\_SONNET\_HAIKU}}}
  & \PH{{{D\_SONNET\_HAIKU}}} \\
GPT-4o vs.\ Haiku
  & \PH{{{CHI2\_GPT4O\_HAIKU}}}
  & \PH{{{P\_GPT4O\_HAIKU}}}
  & \PH{{{D\_GPT4O\_HAIKU}}} \\
\bottomrule
\end{tabular}
\end{table}

\subsection{Figures}

\begin{figure}[ht]
\centering
% \includegraphics[width=0.8\textwidth]{figures/precision_recall_bar.pdf}
\PH{{{FIGURE\_1\_PRECISION\_RECALL\_BAR}}}
\caption{Per-model precision and recall.  Error bars show 95\% bootstrap
  confidence intervals.}
\label{fig:precision-recall}
\end{figure}

\begin{figure}[ht]
\centering
% \includegraphics[width=0.8\textwidth]{figures/hallucination_rates.pdf}
\PH{{{FIGURE\_2\_HALLUCINATION\_RATES}}}
\caption{Hallucination rates per model with 95\% binomial confidence
  intervals.  The dashed line marks the 0.10 hypothesis threshold.}
\label{fig:hallucination}
\end{figure}

\begin{figure}[ht]
\centering
% \includegraphics[width=0.8\textwidth]{figures/stability.pdf}
\PH{{{FIGURE\_3\_STABILITY}}}
\caption{Extraction stability across \PH{{{N\_RUNS}}} independent runs per
  model, measured as the Jaccard similarity of extracted claim sets between
  runs.}
\label{fig:stability}
\end{figure}

% ==========================================================================
\section{Discussion}
\label{sec:discussion}

\subsection{Hypothesis Evaluation}

We pre-registered three hypotheses for the best-performing model:

\begin{enumerate}
  \item \textbf{H1 (Precision $\geq 0.80$):}
    \PH{{{H1\_RESULT}}}---the observed precision is
    \PH{{{PRECISION\_BEST}}} (95\% CI \PH{{{CI\_PREC\_BEST}}}).
  \item \textbf{H2 (Recall $\geq 0.70$):}
    \PH{{{H2\_RESULT}}}---the observed recall is
    \PH{{{RECALL\_BEST}}} (95\% CI \PH{{{CI\_REC\_BEST}}}).
  \item \textbf{H3 (Hallucination rate $\leq 0.10$):}
    \PH{{{H3\_RESULT}}}---the observed hallucination rate is
    \PH{{{HALLUCINATION\_BEST}}} (95\% binomial CI
    \PH{{{CI\_HALL\_BEST}}}).
\end{enumerate}

\subsection{Limitations}

Several limitations qualify these results:

\begin{itemize}
  \item \textbf{Abstract-only text.}  We extract claims from abstracts
    rather than full papers.  Abstracts are information-dense but omit
    methodological nuance; performance may differ on full-text input.
  \item \textbf{LLM-based ground truth.}  Our ground truth is produced by
    LLM annotators rather than domain-expert humans.  While the
    dual-annotator reconciliation protocol mitigates individual model bias,
    systematic blind spots shared across LLMs may persist.
  \item \textbf{Embedding-based matching.}  Cosine similarity at threshold
    0.85 may miss semantically equivalent claims that differ in phrasing
    beyond what the embedding model captures.  A Jaccard token-overlap
    baseline could complement this measure.
\end{itemize}

\subsection{Implications for Downstream Experiments}

\PH{{{IMPLICATIONS\_TEXT}}}

Based on these results, subsequent experiments in the autonomous research
pipeline (E02 and beyond) can adopt \PH{{{BEST\_MODEL}}} as the default
claim-extraction backbone, with an expected precision floor of
\PH{{{PRECISION\_BEST}}} and recall floor of \PH{{{RECALL\_BEST}}}.

% ==========================================================================
\section{Conclusion}
\label{sec:conclusion}

We presented a controlled evaluation of three frontier LLMs on the task of
extracting factual claims from ML paper abstracts.  Using a dual-annotator
ground-truth corpus and embedding-based matching, we measured precision,
recall, F1, and hallucination rates with bootstrap confidence intervals
and pairwise statistical tests.

\PH{{{BEST\_MODEL}}} achieved the highest overall F1 of
\PH{{{F1\_BEST}}}, with a hallucination rate of
\PH{{{HALLUCINATION\_BEST}}}.  \PH{{{CONCLUSION\_EXTRA\_SENTENCE}}}

These baselines inform the design of subsequent modules in the autonomous
research pipeline, including claim verification (E02), knowledge-graph
construction, and hypothesis generation.

% ==========================================================================
\bibliographystyle{plainnat}

\begin{thebibliography}{6}

\bibitem[Zheng et~al.(2023)]{zheng2023judging}
Lianmin Zheng, Wei-Lin Chiang, Ying Sheng, Siyuan Zhuang, Zhanghao Wu,
Yonghao Zhuang, Zi~Lin, Zhuohan Li, Dacheng Li, Eric~P. Xing, Hao Zhang,
Joseph~E. Gonzalez, and Ion Stoica.
\newblock Judging {LLM}-as-a-{J}udge with {MT-B}ench and {C}hatbot {A}rena.
\newblock In \emph{Advances in Neural Information Processing Systems (NeurIPS)},
  2023.

\bibitem[Lu et~al.(2024)]{lu2024aiscientist}
Chris Lu, Cong Lu, Robert~Tjarko Lange, Jakob Foerster, Jeff Clune, and
David Ha.
\newblock The {AI} {S}cientist: Towards fully automated open-ended scientific
  discovery.
\newblock \emph{arXiv preprint arXiv:2408.06292}, 2024.

\bibitem[Beltagy et~al.(2019)]{beltagy2019scibert}
Iz~Beltagy, Kyle Lo, and Arman Cohan.
\newblock {SciBERT}: A pretrained language model for scientific text.
\newblock In \emph{Proceedings of EMNLP}, 2019.

\bibitem[Wadden et~al.(2020)]{wadden2020fact}
David Wadden, Shanchuan Lin, Kyle Lo, Lucy~Lu Wang, Madeleine van Zuylen,
Arman Cohan, and Hannaneh Hajishirzi.
\newblock Fact or fiction: Verifying scientific claims.
\newblock In \emph{Proceedings of EMNLP}, 2020.

\bibitem[McNemar(1947)]{mcnemar1947note}
Quinn McNemar.
\newblock Note on the sampling error of the difference between correlated
  proportions or percentages.
\newblock \emph{Psychometrika}, 12(2):153--157, 1947.

\bibitem[Efron and Tibshirani(1993)]{efron1993bootstrap}
Bradley Efron and Robert~J. Tibshirani.
\newblock \emph{An Introduction to the Bootstrap}.
\newblock Chapman \& Hall, 1993.

\end{thebibliography}

\end{document}
